\begin{quadro}[htbp!]
\centering
\caption{Decisões de adjudicação da auditoria do \textit{dataset}}
\label{quadro:adjudicacao}
\footnotesize
\begin{tabular}{|p{2.3cm}|p{4.9cm}|p{1.5cm}|p{5.3cm}|}
\hline
\textbf{Consultas} & \textbf{Divergência} & \textbf{Decisão} & \textbf{Justificativa} \\
\hline
N36, N37, GF001, GF002, GF003, GF004, GF005, GF006\newline{\scriptsize\itshape anotação independente} & o gabarito classificava a consulta como observacional; o anotador a tratou como fora do domínio, com gabarito determinado 'não chamar' & corrigido & P4 e §3 do manual: consulta sem intenção de busca no acervo brasileiro tem comportamento correto determinado (não chamar); P6 reserva 'observacional' às consultas cujo comportamento correto não é determinável. A consulta passa às métricas principais, na categoria F (fora do domínio). \\ \hline
N33\newline{\scriptsize\itshape anotação independente} & o gabarito exigia project = Mapeamento Sistemático; o anotador anotou só state = São Paulo e listou project como leitura alternativa, pois 'cartografia sistemática' não é nome nem apelido informado na ferramenta & ampliado & As duas leituras são razoáveis (P3): a expressão designa o programa pelo termo distintivo do nome, mas não consta da descrição da ferramenta. project passa a campo opcional; a regra de project do manual passa a registrar o caso. \\ \hline
GF007\newline{\scriptsize\itshape anotação independente} & ambos classificam a consulta como observacional; o gabarito indica 'não chamar' e o anotador, 'chamar sem parâmetros' & mantido & Consulta observacional (P6, valor impossível de representar): o comportamento é reportado, mas não pontuado, de modo que a divergência não afeta nenhuma métrica. Mantém-se a indicação 'não chamar', coerente com N38 (escala 1:10.000.000). \\ \hline
P14\newline{\scriptsize\itshape consistência} & 'detalhado' é anotado como 1:25.000 em todo o dataset, mas P14 aceita 1:25.000 ou 1:50.000 & mantido & As escalas explícitas '(25k ou 50k)' prevalecem sobre o qualificativo 'detalhado' (regra de scale: duas escalas explícitas, aceita-se qualquer uma). \\ \hline
P13\newline{\scriptsize\itshape anotação independente (medida estrita)} & leitura preferencial diferente dentro do conjunto aceito: scale 1:25.000 no gabarito, 1:50.000 no anotador (ambas aceitas por 'maior que 100k') & mantido & Na camada P, a anotação original da equipe do 1º CGEO é mantida como leitura preferencial (manual, §4); a escolha da preferencial não afeta a pontuação. \\ \hline
P21, P22\newline{\scriptsize\itshape anotação independente (medida estrita)} & leitura preferencial diferente dentro do conjunto aceito: sortField creationDate no gabarito, publicationDate no anotador (ambos aceitos, sem pista ligada à ordenação) & mantido & Na camada P, a anotação original da equipe do 1º CGEO é mantida como leitura preferencial (manual, §4); a escolha da preferencial não afeta a pontuação. \\ \hline
GO016\newline{\scriptsize\itshape anotação independente (medida estrita)} & o gabarito aceita sortField publicationDate ou creationDate; o anotador aceitou só publicationDate & mantido & Em 'as cartas mais recentes publicadas entre 2022 e 2023', o verbo qualifica o filtro de período, não a ordenação; sem pista ligada à ordenação, aceitam-se os dois campos (regra de ordenação). É o mesmo raciocínio do anotador em P22. \\ \hline
\end{tabular}
\fonte{Elaborado pelo autor a partir de \texttt{data/auditoria/adjudicacao.json}.}
\end{quadro}
